\section{関連研究}
\label{sec:related_works}

%% ====== 2.1 頻出パターンマイニング ======
\subsection{頻出パターンマイニング}

アイテムセットマイニングの分野では，データ全体での頻度（グローバルサポート）に
基づく頻出パターン抽出が広く研究されてきた~\cite{fournier2017survey}．
代表的な手法として，Apriori~\cite{agrawal1994fast} は反単調性（候補の上位集合の
サポートはその部分集合以下）を利用した候補枝刈りによって探索空間を削減する．
FP-Growth~\cite{han2000mining} は FP-tree と呼ばれる木構造で全パターンを
圧縮表現し，候補生成を不要にする．ECLAT~\cite{zaki1997new} は
縦型データ形式（tidset）を用いて集合積演算でサポートを高速計算する．

これらの手法はグローバルサポートのみを評価するため，
パターンが特定期間に集中して出現する「密集」という現象を捉えることができない．
また，マイニング対象はトランザクション単位のアイテムセットであり，
バスケット（1 回の購買イベント）という粒度を区別する設計にはなっていない．

%% ====== 2.2 時間的・局所的密集パターンマイニング ======
\subsection{時間的・局所的パターンマイニング}

グローバル頻度では捉えられない局所的な頻度変動を検出するための
研究が多数提案されている~\cite{radhakrishna2015survey}．

\textbf{ギャップに基づく手法}:
Periodic Frequent Pattern Mining (PFPM)~\cite{tanbeer2009discovering} は
全体的に高頻度かつ連続出現間のギャップがしきい値以下のパターンを抽出するが，
ギャップしきい値が密集の定義に強く影響するため，パラメータ感度が高い．
Partial Periodic Frequent Pattern Mining (PPFPM)~\cite{kiran2017discovering} は
一部のギャップについてしきい値制約を緩和するが，真に密集した区間を
確実に分離できない場合がある．
LPFIM~\cite{mahanta2005finding} と LPPM~\cite{fournier2021mining} は
連続するギャップがしきい値以下の「ラン」を密集区間として定義するが，
ギャップパラメータの設定に対して精度が不安定という問題がある．

\textbf{スライディングウィンドウに基づく手法}:
先行研究 Apriori-window~\cite{dsaa2025} は，固定幅ウィンドウを
時系列方向にスライドさせながら局所出現頻度を評価する厳密解法であり，
ギャップパラメータを必要とせずに密集パターンを完全列挙する．
本稿はこの手法を直接の拡張基盤として位置づける（\ref{sec:proposed_method}節参照）．

%% ====== 2.3 バスケット構造・マーケットバスケット分析 ======
\subsection{バスケット構造・マルチ粒度トランザクション}
\label{sec:related_basket}

%% --- ここは情報収集が完全ではないため大まかな記述 ---

マーケットバスケット分析（MBA）は，小売業における顧客の購買行動を
アイテムセットとして分析する代表的な応用である~\cite{agrawal1994fast}．
古典的な MBA では「1 顧客 1 購買 = 1 トランザクション」という対応を前提とするが，
実際の POS データやECログでは複数の購買イベントが時刻や日付などの単位で
集約されて提供されることが多い．

アイテムセット間の時系列的な関係を扱う \emph{inter-transaction association rule}
の研究~\cite{lu2000inter} では，複数の隣接トランザクションを一つのウィンドウとして
扱うことで時間をまたいだパターンを抽出する．しかし，この枠組みにおいても
個々のトランザクションは単一のアイテムセットとして扱われており，
1 つのトランザクション内に複数の購買イベントが混在する状況は想定されていない．

\textbf{本研究との差異}:
既存のバスケット分析・パターンマイニング研究の多くは，
(1) 1 トランザクション = 1 アイテムセットという仮定か，
または (2) 複数トランザクションにまたがる時間的関係を扱う枠組みのいずれかを採用している．
これに対して本稿は，\emph{1 トランザクション内に複数の独立したバスケットが存在する}
という現実的なデータ構造を明示的にモデル化し，
そこから生じる偽共起という問題を初めて定式化する点が新規である．

%% ====== 2.4 本研究の位置づけ ======
\subsection{本研究の位置づけ}

表~\ref{tab:related_comparison}に本研究と関連手法の比較を示す．
既存手法がバスケット構造を無視するか，または扱う場合でも偽共起の排除を
目的としていないのに対して，提案手法は両者を統合することで偽共起ゼロを
保証するパターンマイニングを実現する．

\begin{table}[t]
\centering
\caption{関連手法との比較}
\label{tab:related_comparison}
\begin{tabular}{lcccc}
\toprule
手法 & 局所密集 & バスケット構造 & 偽共起排除 & 厳密解 \\
\midrule
Apriori~\cite{agrawal1994fast}           & ✗ & ✗ & ✗ & ✓ \\
FP-Growth~\cite{han2000mining}           & ✗ & ✗ & ✗ & ✓ \\
LPFIM~\cite{mahanta2005finding}          & ✓ & ✗ & ✗ & △ \\
LPPM~\cite{fournier2021mining}           & ✓ & ✗ & ✗ & △ \\
Apriori-window~\cite{dsaa2025}           & ✓ & ✗ & ✗ & ✓ \\
\textbf{Phase 1（提案）}                 & \textbf{✓} & \textbf{✓} & \textbf{✓} & \textbf{✓} \\
\bottomrule
\end{tabular}
\end{table}
